
\subsection{一般Log型の延長ファイバー}
\label{subsec:quadric-log-fibres}
\begin{theorem}[一般対数型の全延長]\label{thm:quadric-generic-log-line}
滑らかな$Q^3$上の$\Log(1,1,1)$成分の非空Zariski開集合では，
規格化された全ambient Poisson延長スキームは被約な$\AA^1$であり，
全延長の一般階数は二である。
\end{theorem}
\begin{proof}
LPTによる一般の対数型は，三独立一次式$\ell_0,\ell_1,\ell_2$と
$\lambda_0+\lambda_1+\lambda_2=0$から
$\alpha=\sum_i\lambda_i\prod_{j\ne i}\ell_j\,d\ell_i$
として得られる\cite{LPT}。
三超平面の共通核$W$への$q$の制限が非退化という非空開条件を課す。
$V=U\oplus W$と$q$について直交分解し，
$U$の座標を$x_0,x_1,x_2$，$W$の座標を$x_3,x_4$と書く。
$q=q_U(x_0,x_1,x_2)+x_3^2+x_4^2$とできる。
$\gamma=d\alpha/3=\iota_\xi(dx_0\wedge dx_1\wedge dx_2)$，
$\xi=ax_0\partial_0+bx_1\partial_1-(a+b)x_2\partial_2$と書く。
$R=x_4\partial_3-x_3\partial_4$，$J=\partial_3\wedge\partial_4$，
$E_U=\sum_{i=0}^2x_i\partial_i$と置けば，
発散零の$q$-Casimir liftと，制限を変えない方向は
\[
 P_0=2\xi\wedge R+\xi(q_U)J,\qquad
 K_R=qJ+E\wedge R=E_U\wedge R+q_UJ
\]
である。従って
\[
 P_t=(2\xi+tE_U)\wedge R+
       \bigl(\xi(q_U)+tq_U\bigr)J.            \tag{L}\label{eq:quadric-log-line}
\]
任意の$U$上の線形場$A$と$U$上の二次式$g$について
$A\wedge R+gJ$はPoissonで，外積二乗も零である。
$q$-Casimir条件は上の式から従う。
よって全$t$で$q$の制限は同一で，一般階数二である。
これらの普遍恒等式は計算コードでも全パラメータで照合する。

この直線が全延長であることを確認する。
$q$と$P_0$を保存する$\mathfrak{so}(q)$のLie代数を考える。
$A\in\mathfrak{so}(q)$が$Q$上の$\sigma=P_0|_Q$を保てば，
$[A,P_0]$は$q$-Casimirかつ発散零で，$Q$上で零となる。
既述の延長核$K_T$と$\delta K_T=-\tfrac32v_T$から$[A,P_0]=0$となる。
従ってintrinsic stabilizerはstrict stabilizerに一致する。
常に$R$が対称性である。一方，
$q=\sum_{i=0}^4x_i^2$，$\xi=(-5x_0\partial_0+4x_1\partial_1+x_2\partial_2)/3$
という正確な点で，$\mathfrak{so}(q)$の十変数に対する
$[A,P_0]=0$の二次係数行列のrankは九であり，
非零九次minorの一つは$10240000/243$である。
行・列は\texttt{check\_log\_quadric.py}の証明書に明示する。
対数源は既約で，$R$は常に存在するので，
一般点のstrict stabilizerはちょうど$\CC R$である。

任意のPoisson延長$P=P_0+K_T$に対して
$[\delta P|_Q,\sigma]=0$がJacobi恒等式の線形帰結として成立する。
$\delta P=-\tfrac32v_T$で，$T\mapsto v_T$は
$\wedge^2V\cong\mathfrak{so}(q)$だから，$K_T=tK_R$に限られる。
この九個の独立な線形条件はJacobi係数イデアルの帰結であり，
全候補の直線上ではJacobi恒等式が恒等的に成立する。
従って延長スキーム自体が被約な$\AA^1$である。
\end{proof}

\begin{corollary}[一般Log延長のambient成分]\label{cor:quadric-log-ambient-membership}
上の$\AA^1$の全点は，$\PP^4$上の
$\Log(1,1,1,2)$型一般階数二族の閉包に属する。
\end{corollary}
\begin{proof}
$A=2\xi+tE_U=\sum_{i=0}^2d_ix_i\partial_i$と置く。
$d=(2a+t,2b+t,-2a-2b+t)$である。
$d$と直交する二つの線形独立な$c=(c_0,c_1,c_2)$に対して
\[
 \vartheta_c=\sum_{i=0}^2c_i\frac{dx_i}{x_i}
 -\frac{c_0+c_1+c_2}{2}\frac{dq}{q}
\]
はEuler水平な閉対数一形式である。
$P_t^\sharp(dx_i/x_i)=d_iR$，$P_t^\sharp(dq)=0$だから
$P_t^\sharp(\vartheta_c)=0$となる。
一般点では二つの$\vartheta_c$が葉の余法線を生成し，
極因子次数は$(1,1,1,2)$となる。
特殊な$t$ではこの一般源の閉包を取る。
一般点の係数に共通の因子がない条件も開であり，
既存rank二Log族への帰属が従う。
\end{proof}

\begin{corollary}[三intrinsic成分の一般延長]\label{cor:quadric-generic-fibre-trichotomy}
滑らかな$Q^3$上のPoisson空間の三既約成分について，
それぞれ非空Zariski開集合における規格化延長スキームは
\[
\begin{array}{c|c|c}
\text{intrinsic成分}&\text{延長スキーム}&\text{一般階数}\\ \hline
\Rat(1,2)&\operatorname{Spec}\CC&2\\
\Log(1,1,1)&\AA^1&\text{全点で }2\\
\Aff&\text{被約二点}&2,\ 4\ (2Q+H)
\end{array}
\]
である。特に一般元にrank四延長が現れるのは$\Aff$成分のみである。
\end{corollary}
\begin{proof}
第一行は系\ref{cor:quadric-generic-rational-rigidity}，第二行は
定理\ref{thm:quadric-generic-log-line}，第三行は
命題\ref{prop:affine-two}とLPTのAffine一般軌道による。
\end{proof}

\subsection{特殊Log境界での交叉二直線}
\label{subsec:quadric-log-boundary-node}
\begin{proposition}[特殊Log境界の二直線延長]\label{prop:quadric-special-log-node}
\[
q_0=\sum_{i=0}^4x_i^2,\quad
\xi=x_0\partial_0+x_1\partial_1-2x_2\partial_2,\quad
R=x_4\partial_3-x_3\partial_4,\quad
S=x_1\partial_0-x_0\partial_1
\]
とし，$J=\partial_3\wedge\partial_4$，
$J_U=\partial_0\wedge\partial_1$，
$P_0=2\xi\wedge R+\xi(x_0^2+x_1^2+x_2^2)J$と置く。
$P_0$が$Q_0=V(q_0)$上に誘導するPoisson構造を$\sigma_0$とすれば，
\[
 \ExtP_{Q_0}(\sigma_0)\cong
 \operatorname{Spec}\CC[t,v]/\bigl(v(t+2)\bigr),
\qquad
 P_{t,v}=P_0+t(q_0J+E\wedge R)+v(q_0J_U+E\wedge S).
\]
これは二本の被約なアフィン直線が$(-2,0)$で横断的に交わるスキームである。
$v=0$の枝は全点で一般階数二，$t=-2,\ v\ne0$の枝は一般階数四で，
その実退化因子は
\[
 D=Q_0+V(x_2)+V(x_2^2+x_3^2+x_4^2).
\]
$\sigma_0$には因子的零点があり一般Log元ではない。
\end{proposition}
\begin{proof}
$\xi$の二つの重みが一致するため，
$S$と$R$は$q_0,P_0$を保つ。
十変数の$A\in\mathfrak{so}(q_0)$に対する
$[A,P_0]=0$の線形行列のrankは八である。
九・十番目の核に対応する二場は$S,R$で，
独立な八次minorは$2304$と正確に計算される。
したがって上述した発散からの必要条件は
全Poisson延長を二変数$P_{t,v}$へ制限する。
この源上でJacobiatorの非零係数は
\[
\begin{aligned}
 J_{013}&=2v(t+2)q_0x_4,&
 J_{014}&=-2v(t+2)q_0x_3,\\
 J_{034}&=2v(t+2)q_0x_1,&
 J_{134}&=-2v(t+2)q_0x_0.
\end{aligned}
\]
これらの係数理想は正確に$(v(t+2))$であり，
スキームに余分な非被約成分はない。
さらに$\iota_{P_{t,v}\wedge P_{t,v}/2}\Omega=f_{t,v}\,dq_0$において
\[
 f_{t,v}=\frac{vx_2}{2}
 (tq_0+2x_0^2+2x_1^2-4x_2^2-4x_3^2-4x_4^2).
\]
従って$v=0$ではPfaffianが零，
$t=-2$では$f=-3vx_2(x_2^2+x_3^2+x_4^2)$となる。
全source・行列minor・Pfaffianを\texttt{check\_log\_quadric.py}で
独立に計算しており，以上は一般Log開集合についての
定理とは矛盾しない。
\end{proof}

\subsection{共鳴対数境界の直交性と全延長ファイバー}
\label{subsec:quadric-log-block-orthogonality}
共鳴する対数型の重み一致と、二つの重み空間の二次形式に関する直交性は別の条件である。
以下では二次形式を六係数のまま動かし、その非退化な配置の全Poisson延長ファイバーを決定する。
直交性は追加のrank四枝の必要十分条件となる。有限個の標本計算を一般定理の根拠とはしない。

\begin{theorem}[共鳴Log型の二次形式に関する直交性判定]
\label{thm:quadric-log-block-orthogonality}
$V=U_1\oplus L\oplus W$を$\dim U_1=\dim W=2$, $\dim L=1$の分解とし、
$U_1$の座標を$u_0,u_1$、$L$の座標を$z$、$W$の座標を$w_0,w_1$とする。
$H\in\operatorname{Sym}^2U_1^*$を非退化な二次形式、
$b\in U_1^*$、$c\in\CC$とし、
\[
 q=q_U+w_0^2+w_1^2,\qquad
 q_U=H(u_0,u_1)+2z\,b(u_0,u_1)+cz^2,
 \qquad\det(q_U)\ne0.
\]
$\xi=u_0\partial_{u_0}+u_1\partial_{u_1}-2z\partial_z$,
$R=w_1\partial_{w_0}-w_0\partial_{w_1}$,
$J=\partial_{w_0}\wedge\partial_{w_1}$,
$E=\sum x_i\partial_{x_i}$とし、
\[
 P_0=2\xi\wedge R+\xi(q_U)J,\qquad
 K_R=qJ+E\wedge R,\qquad Q=V(q)\subset\PP^4.
\]
$P_0$が滑らかな$Q$上に誘導するPoisson構造を$\sigma$とする。
このとき、次がスキームとして成立する。
\[
 \boxed{\quad
 \ExtP_Q(\sigma)\simeq
 \begin{cases}
 \AA^1,&b\ne0,\\[1mm]
 \operatorname{Spec}\CC[t,v]/\bigl(v(t+2)\bigr),&b=0.
 \end{cases}\quad}
\]
$b\ne0$での全延長は$P_0+tK_R$であり、すべて一般階数二である。
$b=0$ではこの階数二の直線に別の階数四の直線が交わり、
後者の一般元の実退化因子は
\[
 Q+V(z)+V(cz^2+w_0^2+w_1^2).
\]
従って、この開いた二次形式の族では、rank四の追加延長が存在するための
必要十分条件は、$U_1$と$L$が$q_U$について直交すること、すなわち$b=0$である。
これは二つの独立な線形条件であり、二次形式全体の上では余次元二の現象となる。
\end{theorem}
\begin{proof}
既述の式\eqref{eq:quadric-log-line}を用いると、全$t$について
\[
 P_t=P_0+tK_R=(2\xi+tE_U)\wedge R+
                  \bigl(\xi(q_U)+tq_U\bigr)J
\]
は$q$をstrict CasimirとするPoisson双ベクトルで、
$P_t\wedge P_t=0$である。しかも$K_R$は$Q$への制限を変えない。

まず$q,P_0$を保つ線形場のLie代数を幾何学的に計算する。
斉次二次の\emph{affine}双ベクトル$P_0$の零スキームの斉次イデアルは
\[
 I_{Z(P_0)}=(u_iw_j\ (i=0,1),\ zw_j\ (j=0,1),\ g),
 \qquad g=\xi(q_U).
\]
実際、$U_1\oplus L$方向と$W$方向の混合係数は
$2\xi\wedge R$の全係数であり、$W$方向の二次係数は$g$である。
$H$の非退化性から$g$は非零二次式である。
この斉次イデアルは
\[
 (u_0,u_1,z)\ \cap\ (w_0,w_1,g)
\]
に等しい。従ってその射影零スキームは互いに交わらない
直線$\PP(W)$と$\PP(U_1\oplus L)$内の次数二の平面曲線$C=V(g)$の和である。
$C$は非被約でもよい。直線と次数二曲線はHilbert多項式が異なり、
さらに$C$は線形方程式を持たずスキーム論的に$\PP(U_1\oplus L)$を張る。
従って$[A,P_0]=0$を満たす$A\in\mathfrak{so}(q)$は
$U_1\oplus L$と$W$をそれぞれ保つ。

その$W$成分は$R$のスカラー倍である。
$U_1\oplus L$成分を$B\in\mathfrak{so}(q_U)$と書くと、
$[B,P_0]=0$の混合係数は$[B,\xi]\wedge R$であるから
$[B,\xi]=0$である。逆にこの二条件は$B(g)=0$を含意するので十分である。
$\xi$の固有値$1,-2$の分解から
$B=\operatorname{diag}(D,d)$と書ける。
$q_U$を保つ条件は正確に
\[
 D^tH+HD=0,\qquad D^tb+db=0,\qquad 2cd=0.
\]
$c\ne0$なら$d=0$であり、$D\in\mathfrak{so}(H)$となる。
$H$は二次元で非退化だから$\mathfrak{so}(H)$は一次元で、
非零元はすべて可逆である。従って$b\ne0$なら$D=0$である。
一方$c=0$なら$q_U$の非退化性から
$b^tH^{-1}b\ne0$となる。このとき$D\in\mathfrak{so}(H)$が非零なら、
$D^tb+db=0$により$b$は$D^t$の非零固有値の固有ベクトルとなるが、
反対称な$D$の非零固有値の固有ベクトルは$H^{-1}$に関して等方的である。
これは$b^tH^{-1}b\ne0$に矛盾する。
よってこの場合も$D=d=0$である。
$b=0$なら$q_U$非退化性から自動的に$c\ne0$で、
$D\in\mathfrak{so}(H)$の一次元が残る。
結局
\[
 \mathfrak{aut}_{\mathrm{strict}}(q,P_0)
 =\begin{cases}\CC R,&b\ne0,\\
  \CC R\oplus\mathfrak{so}(H),&b=0.
 \end{cases}
\]

このstrict stabilizerは$Q$上の$\sigma$の
$\mathfrak{so}(q)$-stabilizerに等しい。
実際、$A$が$\sigma$を保存すれば$[A,P_0]$は発散零で$q$-Casimirかつ
$Q$への制限が零である。既述の核$K_T$と
$\delta K_T=-\frac32v_T$の単射性から$[A,P_0]=0$を得る。
任意のambient Poisson延長について、発散を取ったJacobiの線形帰結により
その発散場は$\sigma$のこのstabilizerに入る。
従って$b\ne0$では全候補が$P_0+tK_R$で尽くされ、
九個の独立な線形条件がJacobiイデアルから出る。
上の直線ではJacobiが恒等零だから、延長スキームは被約な$\AA^1$である。

$b=0$ならば$H$を二次元の変数変換で$u_0^2+u_1^2$に、
$cz^2$を$z^2$に直すことができる。これらの変数変換は$\xi$と可換であり、
$P_0$と制限・延長ファイバーを先の標準模型へ移す。
命題\ref{prop:quadric-special-log-node}の全Jacobi理想
$(v(t+2))$を適用すれば、二直線交叉とrank四の因子の式を得る。
\end{proof}

\begin{corollary}[一変数の直交性wall]\label{cor:quadric-resonant-log-wall}
一変数族$q_c=\sum x_i^2+2cx_0x_2$は$H=u_0^2+u_1^2$、
$b(u)=cu_0$、$z=x_2$、二次係数$1$の場合であり、
この定理の一次元切断に当たる。
$c^2\ne1$なら
\[
\ExtP_{Q_c}(\sigma_c)\cong
\begin{cases}\AA^1,&c\ne0,\\
\operatorname{Spec}\CC[t,v]/(v(t+2)),&c=0.
\end{cases}
\]
前稿のexact九次minorは$13824c/(c^2-1)^4$であり、
この結論の独立計算上のチェックとして保持できる。
\end{corollary}
\begin{remark}
より一般の二次形式の族では直交性条件$b=0$は余次元二である。
$H$の退化を許す層、または$q_U$が退化する層は本定理の射程外である。
$c=0$自体は除外しない。$q_U$が非退化なら$c=0$でも$b$は非等方となり、
追加のstrict対称性は存在しない。
今回の独立計算は五つの正確な二次形式についてstrict stabilizerの階数と全$P_t$の
Jacobi/Pfaffianを検査したもので、一般定理の証明は上記の零スキームの論法による。
新規性の一次文献調査は未完了である。
\end{remark}


\subsection{滑らかな二次超曲面における等方的なLog境界と非被約延長}
\label{subsec:quadric-log-isotropic}
先の共鳴Log模型で$q=q_U+w_0^2+w_1^2$と直交分解している限り，
$q_U$の退化は$q$そのものの退化を意味する。従って滑らかな
$Q=V(q)$の新しい境界は，対数構造の三線形形式
$(u_0,u_1,z)$の共通核$W=\operatorname{Span}(w_0,w_1)$に
\emph{制限した}二次形式$q|_W$の退化として調べる必要がある。
この層では$Q$が滑らかなまま，延長ファイバーの既約成分だけでなく
scheme構造まで跳躍することを次に示す。

$V=\CC^5$の座標を$(u_0,u_1,z,w_0,w_1)$とし，
$\xi=u_0\partial_{u_0}+u_1\partial_{u_1}-2z\partial_z$，
$\gamma=\iota_\xi(du_0\wedge du_1\wedge dz)$と置く。
非退化二次式$q$ごとに
\begin{equation}
 \iota_{P_0(q)}(du_0\wedge du_1\wedge dz\wedge dw_0\wedge dw_1)
  =dq\wedge\gamma                                         \label{eq:isotropic-log-p0}
\end{equation}
と定める。$P_0(q)$は$q$をstrict Casimirとする
発散零の階数二Poisson双ベクトルであり，$Q=V(q)$に
非零の共鳴対数型Poisson構造$\sigma_q$を誘導する。

\begin{theorem}[制限二次形式の退化に伴う三種類の延長ファイバー]
\label{thm:quadric-isotropic-log-fibres}
以下の三つの非退化二次式について，固定された
$\sigma_q=P_0(q)|_Q$の\emph{全}ambient Poisson延長スキームは
次の表の通りである。すべて$Q$は滑らかであり，
$W$は同一の二次元部分空間である。
\[
\begin{array}{c|l|c|c}
 &q&\operatorname{rank}(q|_W)&\ExtP_Q(\sigma_q)\\ \hline
 \mathrm{I}&u_0^2+u_1^2+z^2+2u_0w_0+w_1^2&1&\AA^2_{a,b}\\
 \mathrm{II}&u_0^2+2u_1w_0+2zw_1+w_0^2&1&
     \operatorname{Spec}\CC[a,b]/(a(b-4))\\
 \mathrm{III}&u_0^2+2u_1w_0+2zw_1&0&
     \operatorname{Spec}\CC[a,b,c]/(ac,c^2)
\end{array}
\]
Iでは$a\ne0$の延長は一般階数四で，その実退化因子は
$Q+2V(u_0)+V(z)$である。IIでは$b=4,\ a\ne0$の枝が
一般階数四で，その実退化因子は
\[
 Q+V(z)+V\bigl((u_1+w_0)^2+2zw_1\bigr)
\]
である。IIIでは被約台は$\AA^2_{a,b}$で，$a\ne0$が一般階数四，
実退化因子は$Q+2V(u_1)+V(z)$である。しかしファイバー自身は
非被約であり，$a=c=0$の直線上に$c$という平方零方向を持つ。
特に$\operatorname{rank}(q|_W)$だけでは
Poisson延長ファイバーのscheme構造を決定できない。
\end{theorem}
\begin{proof}
$\gamma$は三次元空間の二形式で，$d\gamma=0$である。
従って\eqref{eq:isotropic-log-p0}は閉じた分解可能三形式を与え，
対応する双ベクトルは$q$をCasimirとし，発散零でPoissonである。
$q$の非退化性は，座標順$(u_0,u_1,z,w_0,w_1)$で
半Hessian行列$G=\operatorname{Hess}(q)/2$の行列式が順に
$-1,1,1$となることから従う。

任意の$q$-Casimir延長は$P_0(q)$との差を，一意に
\[
 K_A=qT_A+\tfrac12 E\wedge(Ax),\qquad
 T_A=\tfrac12 A G^{-1},\qquad A^tG+GA=0
 \tag{K}\label{eq:isotropic-K}
\]
と書ける。ここでは$T_A$を定数反対称双ベクトルと見なし，
$E$はEuler場である。$K_A$の$q$への制限は零で，
$\delta K_A=\tfrac32 Ax$となる。これは前小節の全延長核と
定数双ベクトルの規約を揃えた書き方である。

$P_0(q)$は発散零なので，延長$P_0(q)+K_A$がPoissonなら，
発散を取ったJacobi恒等式と$Q$への制限から
$[Ax|_Q,\sigma_q]=0$が必要である。
逆に$q$を保つ線形場$Ax$が$\sigma_q$を保てば，
$[Ax,P_0(q)]$は$q$-Casimirかつ発散零で，$Q$上で零となる。
前小節の核の発散写像は単射だから，$[Ax,P_0(q)]=0$である。
従って\emph{全}Poisson延長の候補空間は
$\mathfrak{so}(q)$内のstrict stabilizerのアフィン空間へ
scheme論的な線形方程式で埋め込まれる。

三例のstrict stabilizerはそれぞれ次の線形場により張られる。
\[
\begin{array}{c|l}
\mathrm I&-u_0\partial_{u_1}+u_1\partial_{w_0},\quad
               w_1\partial_{w_0}-u_0\partial_{w_1}\\
\mathrm{II}&u_1\partial_{u_0}+u_0\partial_{u_1}-u_0\partial_{w_0},\quad
               -z\partial_{w_0}+(u_1+w_0)\partial_{w_1}\\
\mathrm{III}&u_1\partial_{u_0}-u_0\partial_{w_0},\quad
               z\partial_{w_0}-u_1\partial_{w_1},\quad
               u_1\partial_{u_1}-z\partial_z-w_0\partial_{w_0}+w_1\partial_{w_1}
\end{array}
\]
実際$A=G^{-1}Z$，$Z$を十変数の反対称行列とし，
$[Ax,P_0(q)]$の全二次係数に対応する線形行列を作ると，
行列サイズと階数は順に$(68\times10,8)$，$(68\times10,8)$，
$(47\times10,7)$である。非零の階数最大小行列式は
順に$2304,-2304,3456$となる。これらの行・列位置と
式は\texttt{check\_isotropic\_log\_boundaries.py}が
前提とする既存producerをimportせず再構成する。

表の各場に\eqref{eq:isotropic-K}を適用し，その順に
パラメータ$a,b$または$a,b,c$を掛けて$P_0$に加える。
得られた候補に対する\emph{全}Jacobi係数の生成イデアルは
\[
 I_{\mathrm I}=(0),\qquad
 I_{\mathrm{II}}=(a(b-4)),\qquad
 I_{\mathrm{III}}=(ac,c^2).
\]
第IIIの場合，例えば$(0,2,3)$成分は
$\frac12 ac\,z\,q$，$(1,2,3)$成分は
$\frac12 c^2 z\,q$であり，残りも
これら二式のイデアルに属する。
従って上記は零集合の比較ではなく，Jacobiイデアルの等式である。

$\beta=\iota_{P\wedge P/2}\Omega=f\,dq$と書くと，
Iでは$f=\frac32a u_0^2z$となる。
IIの第二枝$b=4$では
$f=\frac32a z((u_1+w_0)^2+2zw_1)$である。
IIIの被約台$c=0$では$f=\frac32a u_1^2z$となる。
$q$のEuler次数は二なので，実Pfaffianは$2qf$である。
これで各階数四の開部分と実退化因子を得る。
IIIでは$(ac,c^2)=(c)\cap(a,c)^2$であり，
被約台の平面に，直線$a=c=0$に沿うembeddedな
平方零方向が付く。すべての等式は有理数係数の恒等式として
付属証明書で再現した。
\end{proof}

\begin{remark}
この結果は特定の三つの\emph{滑らかな}二次超曲面と
同一の共鳴Log構成に関する全延長ファイバーの分類であり，
$q|_W$が退化する全ての滑らかなquadricの分類ではない。
特にIIIの非被約方向は新しい被約既約成分を意味しない。
この現象は，これまでの147成分という\emph{被約}分類と，
Poissonスキームのembedded構造の違いを示す。
原稿の通常の論証と有限の正確係数計算を区別し，文献上の新規性は未確認とする。
\end{remark}


\subsection{非被約延長ファイバーからのambient Poisson二次障害}
\label{subsec:quadric-ambient-second-order-obstruction}
前節の完全等方Log例では，固定されたquadricとその上の
Poisson構造への延長ファイバーが
$\operatorname{Spec}\CC[a,b,c]/(ac,c^2)$となった。
$c$の平方零方向について，この二次障害が固定制限を外しても
残ることを示す。これは$\PP^4$上の\emph{全}Poisson変形の問題であり，
単なる相対延長イデアルの再計算ではない。

\begin{theorem}[ambientでも消えない二次障害]
\label{thm:isotropic-ambient-second-order-obstruction}
座標を$X=(u_0,u_1,z,w_0,w_1)$とし，
\[
 q=u_0^2+2u_1w_0+2zw_1,\qquad
 \xi=u_0\partial_{u_0}+u_1\partial_{u_1}-2z\partial_z,
\]
\[
 \iota_{P_0}\Omega=dq\wedge
 \iota_\xi(du_0\wedge du_1\wedge dz)
\]
とする。$E=\sum x_i\partial_{x_i}$，$G=\operatorname{Hess}(q)/2$，
\[
 C=u_1\partial_{u_1}-z\partial_z-w_0\partial_{w_0}
 +w_1\partial_{w_1},\qquad
 A=\operatorname{diag}(0,1,-1,-1,1),
\]
\[
 T=\tfrac12 AG^{-1},\qquad
 K=qT+\tfrac12 E\wedge C
\]
と置く。$[P_0]$を$\PP^4$のPoissonスキーム$\mathcal P$の点とすると，
次が成り立つ。
\begin{enumerate}
\item $K$は$\mathcal P$の$[P_0]$における非零の一次Poisson変形を与える。
これは前節の延長ファイバーのembeddedな$c$方向である。
\item 任意の斉次二次双ベクトル$R$について
\[
 P_0+\varepsilon K+\varepsilon^2R
 \quad\pmod{\varepsilon^3}
\]
が$\PP^4$上のPoisson構造になることはない。
従って$c$方向は，制限quadricや誘導Poisson構造を変えることまで
許しても二次まで延長できない。
\item 特に$\mathcal P$は$[P_0]$で滑らかではない。
全150個の斉次二次双ベクトル係数から計算した
射影Jacobi線形化の階数は$95$であり，
$T_{[P_0]}\mathcal P$の次元は$30$である。
\end{enumerate}
\end{theorem}
\begin{proof}
二次式$q$は非退化で，$A^tG+GA=0$となる。
従って$T$は反対称で，$K$は$q$をstrict Casimirとする
双ベクトルである。前節の全延長の式から，$K$は零でない
$c$方向を表し，$P_0+\varepsilon K$は
$\CC[\varepsilon]/(\varepsilon^2)$上のPoisson構造となる。

一般の斉次二次双ベクトル$B$に対して
\[
 J(B)_{ijk}=\sum_{\ell=0}^4
 \bigl(B_{i\ell}\partial_\ell B_{jk}
 +B_{j\ell}\partial_\ell B_{ki}
 +B_{k\ell}\partial_\ell B_{ij}\bigr)
\]
と定義する。この符号規約では$\PP^4$上の射影Jacobi方程式は
$\mathcal F(B)=E\wedge J(B)=0$である。
$J(P_0)=0$および$J(P_0+K)=J(K)$は直接確認できる。
従って二次補正$R$の存在は
\begin{equation}
 E\wedge\bigl(dJ_{P_0}(R)+J(K)\bigr)=0
 \label{eq:isotropic-ambient-second-obstruction}
\end{equation}
を必要とする。ここで$R$は\emph{固定quadricや固定制限の条件を
一切課していない}任意の150係数の双ベクトルである。

\noindent
次数四の四重ベクトル$F$の成分を$F_{ijkl}$と書き，
$[m]F_{ijkl}$を単項式$m$の係数とする。次の五係数汎関数を取る。
\[
\begin{split}
 \Lambda(F)={}&2[u_0^3z]F_{0123}
 +[u_0u_1zw_0]F_{0123}+[u_0z^2w_1]F_{0123}\\
 &+4[u_0^3u_1]F_{0124}
 +10[u_0u_1^2w_0]F_{0124}.
\end{split}
\]
任意の$R$に対し
\begin{equation}
 \Lambda\bigl(E\wedge dJ_{P_0}(R)\bigr)=0
 \label{eq:isotropic-five-term-cancellation}
\end{equation}
となる。実際，150個の基底
$x_ax_b\partial_i\wedge\partial_j$上で計算すると，
この五係数に現れる非零列は六列だけで，係数行列は
\[
\begin{pmatrix}
 2&0&-12&0&2&0\\
 -2&12&0&-12&4&-6\\
 -2&0&0&0&-8&6\\
 0&2&6&-2&0&0\\
 0&-2&0&2&0&0
\end{pmatrix}.
\]
その左から$(2,1,1,4,10)$を掛けると零になる。
これは有限標本ではなく，\emph{全150係数}に対する係数恒等式である。
一方，$J(K)$の非零成分は
\[
\begin{aligned}
 J(K)_{123}&=\tfrac12zq,&
 J(K)_{124}&=-\tfrac12u_1q,\\
 J(K)_{134}&=\tfrac12w_1q,&
 J(K)_{234}&=-\tfrac12w_0q.
\end{aligned}
\]
従って，上の五係数に対する$E\wedge J(K)$の値は順に
$(\frac12,1,1,-\frac12,-1)$であり，
\[
 \Lambda\bigl(E\wedge J(K)\bigr)
 =2\cdot\tfrac12+1+1+4\cdot(-\tfrac12)+10\cdot(-1)=-9\ne0.
\]
式\eqref{eq:isotropic-ambient-second-obstruction}は
\eqref{eq:isotropic-five-term-cancellation}と矛盾する。
これで第2項が示された。

全350個の射影Jacobi係数を150個の斉次二次係数で線形化すると
階数$95$となる。$E\wedge$一次ベクトル場の24次元のゲージと
射影尺度の一次元を引けば，接空間次元は$150-95-24-1=30$である。
この係数階数と五係数恒等式は，付属する独立の
\texttt{check\_ambient\_quadratic\_obstruction.py}で$\mathbb Q$上に再構成できる。
滑らかなスキームでは任意の一次変形が
$\CC[\varepsilon]/(\varepsilon^3)$まで持ち上がるので，
上の非持上げ性は$[P_0]$の非滑らかさを含意する。
\end{proof}

\begin{remark}[非被約性と特異性の区別]
\label{rem:ambient-obstruction-not-nonreduced}
本定理はambient Poissonスキームの\emph{二次障害が非零}であることを
証明するが，ambientの局所環が非被約であることは証明しない。
被約な特異スキームにも二次まで持ち上がらない接方向はある。
固定quadricの延長ファイバーの非被約性は前節の別の定理であり，
両者を混同しない。なお，被約台上の一般rank四点が
既存147成分のどれに属するかは，ここでは決定していない。
\end{remark}


\begin{corollary}[完全等方Log点での二次接錐と局所次元]
\label{cor:ambient-isotropic-tangent-cone-bound}
定理\ref{thm:isotropic-ambient-second-order-obstruction}の
$p=[P_0]\in\mathcal P=\mathbb P\operatorname{Pois}(\PP^4)$に対し、
既述の五係数汎関数$\Lambda$から定まる二次式
\[
 \mathcal Q_\Lambda(B)=\Lambda\bigl(E\wedge J(B)\bigr)
\]
は$30$次元のZariski接空間$T_p\mathcal P$上の非零二次形式を定め、
$p$での接錐はその零集合に含まれる。従って
\[
 \boxed{\dim T_p\mathcal P=30,\qquad\dim_p\mathcal P\le29.}
\]
特に$p$を含む既約成分の次元はすべて$29$以下である。
現行カタログの次元$36$の\texttt{R2\_RAT113}と
次元$39$の\texttt{R2\_LPB}は$p$を含まない。

さらに完全等方Logの固定$Q$延長ファイバーを表す場
\[
 A_1=u_1\partial_{u_0}-u_0\partial_{w_0},\quad
 A_2=z\partial_{w_0}-u_1\partial_{w_1},\quad A_3=C
\]
に対応する式\eqref{eq:isotropic-K}の核双ベクトルを
$K_1,K_2,K_3$とすると、
\begin{equation}\label{eq:isotropic-tangent-slice}
 \mathcal Q_\Lambda(P_0+aK_1+bK_2+cK_3)=-9c^2
\end{equation}
が$a,b,c$の恒等式として成立する。
\end{corollary}
\begin{proof}
全$150$係数で成立する五係数消滅恒等式
\eqref{eq:isotropic-five-term-cancellation}は
$\Lambda(d\mathcal F_{P_0}(B))=0$を意味する。
$\mathcal F$は二次で$\mathcal F(P_0)=0$なので、
局所座標$P_0+B$では
$\Lambda(\mathcal F(P_0+B))=\mathcal Q_\Lambda(B)$となる。

任意の線形ベクトル場$V$に対し、
$[E,B]=[E,V]=0$とSchouten括弧のLeibniz則から
$E\wedge J(B+E\wedge V)=E\wedge J(B)$である。
また上述の線形化の消滅から
$\mathcal Q_\Lambda(B+\lambda P_0)=\mathcal Q_\Lambda(B)$となる。
従ってこの二次式はEulerゲージおよび射影尺度の商である
$T_p\mathcal P$に降りる。
前定理の接方向$K=K_3$で
$\mathcal Q_\Lambda(K)=-9\ne0$なので、
接錐の定義イデアルには接空間上の非零二次式が含まれる。
局所環と接錐のKrull次元は同じであるから
$\dim_p\mathcal P\le30-1=29$を得る。
カタログの二成分は次元がこの上界を超えるので$p$を含まない。

最後に$A_1,A_2,A_3$を式\eqref{eq:isotropic-K}に代入し、
$E\wedge J(P_0+aK_1+bK_2+cK_3)$の
$\Lambda$で抽出する五係数を直接展開すると
\[
 (c^2/2,\ c^2,\ c^2,\ -c^2/2,\ -c^2)
\]
となる。これに重み$(2,1,1,4,10)$を掛ければ$-9c^2$。
この恒等式は元の障害producerとは独立にSymPyで再計算した。
\end{proof}
\begin{remark}[未解決の局所幾何]
この系は局所次元の等号成立、接錐の既約分解、
$p$を含む残る成分の特定、ambient局所環の非被約性を意味しない。
表示した二次式が検出するのは固定ファイバーの$c^2$関係であり、
$ac$関係のambient障害性も未確認である。
\end{remark}

